%%%%%%%%%%%%%%%%%%%%%%%%%%%%%%%%%%%%%%%%%%%%%%%%%%%%%%%%%%%%%%%%%%%%%
%  RISE_main_kr.tex  —  Korean 번역본 (한글본, standalone)
%  RISE_main.tex(영문본)의 한글 번역본. 영문본이 바뀌면 이 파일도 함께 갱신한다.
%
%  Springer sn-jnl 템플릿은 한글을 렌더링하지 못하므로(pdflatex 는 CJK 미지원,
%  클래스가 LuaLaTeX 비호환) 한글본은 article 클래스 + LuaLaTeX + BibTeX 로
%  독립 구성한다. 참고문헌은 영문본과 "동일한" sub/references-sn.bib 를 사용한다.
%  The Springer sn-jnl class cannot render Korean, so the Korean version is a
%  standalone article-class file built with LuaLaTeX + BibTeX (same .bib database).
%
%  Build / 빌드:  lualatex -> bibtex -> lualatex x2
%  (제공된 RISE_main_build.sh / .cmd 의 `kr` 명령이 자동 처리)
%%%%%%%%%%%%%%%%%%%%%%%%%%%%%%%%%%%%%%%%%%%%%%%%%%%%%%%%%%%%%%%%%%%%%
\documentclass[11pt,a4paper]{article}

\usepackage[a4paper,margin=25mm]{geometry}
\usepackage{setspace}\onehalfspacing

%% Korean fonts (LuaLaTeX/fontspec). Fall back to fonts guaranteed on Windows.
%% 한글 글꼴(LuaLaTeX/fontspec). 윈도우 기본 글꼴로 안전하게 대체한다.
\usepackage{fontspec}
\defaultfontfeatures{Ligatures=TeX}
\IfFontExistsTF{Noto Serif KR}{\setmainfont{Noto Serif KR}}{%
  \IfFontExistsTF{NanumMyeongjo}{\setmainfont{NanumMyeongjo}}{%
    \IfFontExistsTF{Batang}{\setmainfont{Batang}}{\setmainfont{Malgun Gothic}}}}
\IfFontExistsTF{Malgun Gothic}{\setsansfont{Malgun Gothic}}{%
  \IfFontExistsTF{Noto Sans KR}{\setsansfont{Noto Sans KR}}{\setsansfont{Gulim}}}
\IfFontExistsTF{Consolas}{\setmonofont{Consolas}}{}

%% Packages used by the body / 본문이 사용하는 패키지
\usepackage{graphicx}
\usepackage{booktabs}
\usepackage{array}
\usepackage{tabularx}
\usepackage{threeparttable}
\usepackage{multirow}
\usepackage{amsmath,amssymb}
\usepackage{xcolor}
\usepackage{tikz}\usetikzlibrary{arrows.meta}
\usepackage{caption}
\usepackage[round]{natbib}% author-year (APA), matching the English manuscript
\usepackage[unicode,hidelinks]{hyperref}
\usepackage{etoolbox}

%% Float labels in Korean / 플로트 이름 한글화
\renewcommand{\figurename}{그림}
\renewcommand{\tablename}{표}
\renewcommand{\refname}{참고문헌}

%% Table helpers mirrored from the English template / 영문본의 표 헬퍼 이식
\newcolumntype{L}{>{\raggedright\arraybackslash}X}
\newcommand{\TableNote}[1]{\par\vspace{2pt}\noindent\parbox{\linewidth}{\footnotesize\raggedright\textit{주.}\ #1}}

%% Unnumbered section heads, matching the English manuscript style.
%% 영문본과 동일하게 절 제목 번호를 제거한다.
\setcounter{secnumdepth}{0}

% =========================================================================
% Front matter — 제목/저자/초록/키워드를 새 원고에 맞게 수정 (일반 placeholder)
% =========================================================================
\title{\textbf{논문 전체 제목을 여기에:\\
부제가 있다면 여기에}}
\author{제1저자\textsuperscript{1} \and 교신저자\textsuperscript{1,$*$} \and 제3저자\textsuperscript{2}}
\date{\normalsize \textsuperscript{1}소속 기관 \quad \textsuperscript{2}제2 소속 기관 \quad
  \textsuperscript{$*$}교신저자: corresponding.author@example.edu}

\begin{document}
\maketitle

\begin{abstract}
\noindent
초록을 하나의 단락으로 작성합니다. 연구가 다루는 문제와 배경, 연구 목적과 연구
문제, 연구 대상과 방법, 핵심 결과, 그리고 과학교육 연구와 실천에 대한 함의를
압축하여 기술합니다. 분량 제한은 저널의 최신 저자 안내에서 확인합니다. 이 한글본은
영문본(RISE\_main.tex)의 번역본이므로, 영문본이 수정되면 이 초록도 함께 갱신합니다.
\par\smallskip
\noindent\textbf{주제어:} 주제어1; 주제어2; 주제어3; 주제어4
\end{abstract}

% =========================================================================
% 본문 — 각 절 파일은 sub/ 폴더에 있습니다(한글 안내만 포함).
% =========================================================================
%%%%%%%%%%%%%%%%%%%%%%%%%%%%%%%%%%%%%%%%%%%%%%%%%%%%%%%%%%%%%%%%%%%%%
%  sub/kr-1-introduction.tex  —  1. 서론 (한글본)
%  영문본 sub/en-1-introduction.tex 의 번역본. 구조를 동일하게 유지한다.
%%%%%%%%%%%%%%%%%%%%%%%%%%%%%%%%%%%%%%%%%%%%%%%%%%%%%%%%%%%%%%%%%%%%%
\section{서론}\label{sec:introduction-kr}

% 서론에서는 연구가 다루는 과학교육 문제와 그 중요성을 제시하고, 선행연구를
% 비판적으로 검토하여 연구 공백(gap)을 드러낸 뒤, 본 연구의 목적과 연구 문제를
% 명확히 밝힙니다.

여기에 서론을 작성합니다. 연구가 다루는 과학교육 문제와 그 중요성을 밝히고,
선행연구를 검토하여 연구 공백을 드러낸 뒤, 명시적인 연구 문제로 마무리합니다.

% 인용은 \citep{key}(괄호 인용) 또는 \citet{key}(본문 인용)로 씁니다.
% 예: 선행연구는 이 문제를 다루었으며 \citep{example2024}, \citet{example2024}는
% \dots 라고 논하였다.
   % 1. 서론
%%%%%%%%%%%%%%%%%%%%%%%%%%%%%%%%%%%%%%%%%%%%%%%%%%%%%%%%%%%%%%%%%%%%%
%  sub/kr-2-framework.tex  —  2. 이론적 배경 (한글본)
%  영문본 sub/en-2-framework.tex 의 번역본.
%%%%%%%%%%%%%%%%%%%%%%%%%%%%%%%%%%%%%%%%%%%%%%%%%%%%%%%%%%%%%%%%%%%%%
\section{이론적 배경}\label{sec:framework-kr}

% 연구를 뒷받침하는 이론적 틀, 핵심 개념, 분석 렌즈를 제시합니다. 연구 문제와
% 방법(코딩 체계, 루브릭 등)을 이 틀에 연결하면 이후 결과 해석의 근거가 됩니다.

연구를 뒷받침하는 이론적 틀, 핵심 개념, 분석 렌즈를 제시하고, 연구 문제와 연구
방법을 이 틀에 연결합니다 \citep{examplebook2020}.
      % 2. 이론적 배경
%%%%%%%%%%%%%%%%%%%%%%%%%%%%%%%%%%%%%%%%%%%%%%%%%%%%%%%%%%%%%%%%%%%%%
%  sub/kr-3-methods.tex  —  3. 연구 방법 (한글본)
%  영문본 sub/en-3-methods.tex 의 번역본.
%%%%%%%%%%%%%%%%%%%%%%%%%%%%%%%%%%%%%%%%%%%%%%%%%%%%%%%%%%%%%%%%%%%%%
\section{연구 방법}\label{sec:methods-kr}

% 연구 설계, 맥락과 참여자, 자료 수집, 분석 절차(코딩 체계·신뢰도·타당도)를
% 재현 가능하도록 기술합니다. 연구윤리(IRB) 승인과 동의 절차도 밝힙니다.

연구 설계, 맥락과 참여자, 자료 원천, 분석 절차(코딩 체계, 신뢰도, 타당도)를
재현 가능한 수준으로 상세히 기술합니다.

% 표 예시 (필요 시 사용):
% \begin{table}[htbp]
%   \centering
%   \caption{참여자 요약}\label{tab:participants-kr}
%   \begin{tabularx}{\linewidth}{@{}lL@{}}
%     \toprule
%     집단 & 설명 \\
%     \midrule
%     A & \dots \\
%     \bottomrule
%   \end{tabularx}
%   \TableNote{필요하면 표 아래에 주를 추가합니다.}
% \end{table}
        % 3. 연구 방법
%%%%%%%%%%%%%%%%%%%%%%%%%%%%%%%%%%%%%%%%%%%%%%%%%%%%%%%%%%%%%%%%%%%%%
%  sub/kr-4-results.tex  —  4. 연구 결과 (한글본)
%  영문본 sub/en-4-results.tex 의 번역본.
%%%%%%%%%%%%%%%%%%%%%%%%%%%%%%%%%%%%%%%%%%%%%%%%%%%%%%%%%%%%%%%%%%%%%
\section{연구 결과}\label{sec:results-kr}

% 연구 문제 순서에 따라 결과를 제시합니다. 표와 그림은 본문에서 먼저 언급한 뒤
% 배치하고, 모든 표/그림에 \label 을 붙여 \ref 로 참조합니다. 해석은 최소화하고
% 근거(수치·인용·사례)를 중심으로 서술합니다.

연구 문제 순서에 따라 결과를 제시합니다. 모든 표와 그림은 본문에서 먼저 언급한
뒤에 배치합니다.

% 그림 예시 (필요 시 사용):
% \begin{figure}[htbp]
%   \centering
%   \includegraphics[width=\linewidth]{images/your-figure}
%   \caption{그림을 설명합니다.}\label{fig:example-kr}
% \end{figure}
        % 4. 연구 결과
%%%%%%%%%%%%%%%%%%%%%%%%%%%%%%%%%%%%%%%%%%%%%%%%%%%%%%%%%%%%%%%%%%%%%
%  sub/kr-5-discussion.tex  —  5. 논의 및 결론 (한글본)
%  영문본 sub/en-5-discussion.tex 의 번역본.
%%%%%%%%%%%%%%%%%%%%%%%%%%%%%%%%%%%%%%%%%%%%%%%%%%%%%%%%%%%%%%%%%%%%%
\section{논의 및 결론}\label{sec:discussion-kr}

% 핵심 결과를 연구 문제·이론 틀과 연결해 해석하고, 선행연구와의 관계(확장/반박),
% 과학교육 연구·실천에의 함의, 연구의 한계와 후속 연구 방향을 제시합니다.

핵심 결과를 연구 문제 및 이론적 틀과 연결하여 해석하고, 선행연구와의 관계를
논의하며 \citep{examplechapter2019}, 과학교육 연구와 실천에 대한 함의, 연구의
한계와 후속 연구 방향을 제시합니다.
     % 5. 논의 및 결론

% =========================================================================
% 연구 윤리 선언 (Declarations) — 원고에 맞게 작성
% =========================================================================
\section{연구 윤리 선언}
\begin{itemize}
  \item \textbf{연구비.} 연구비 출처를 밝히거나, 외부 연구비를 받지 않았음을 기술합니다.
  \item \textbf{이해상충.} 이해상충이 있으면 밝히고, 없으면 없음을 기술합니다.
  \item \textbf{연구윤리 승인.} IRB/연구윤리 승인 또는 면제 정보를 기술합니다.
  \item \textbf{연구참여 동의.} 연구참여 동의(미성년자는 보호자 동의 포함) 취득을 기술합니다.
  \item \textbf{데이터 공개.} 데이터 접근 또는 요청 방법을 기술합니다.
  \item \textbf{생성형 AI 사용.} 생성형 AI의 사용 여부와 범위를 밝힙니다.
\end{itemize}

%% References / 참고문헌 (영문본과 동일한 데이터베이스 사용).
\bibliographystyle{apalike}% author-year (APA-style), matching the English manuscript
\bibliography{sub/references-sn}

\end{document}
%%%%%%%%%%%%%%%%%%%%%%%%%%%%%%%%%%%%%%%%%%%%%%%%%%%%%%%%%%%%%%%%%%%%%
%  END RISE_main_kr.tex
%%%%%%%%%%%%%%%%%%%%%%%%%%%%%%%%%%%%%%%%%%%%%%%%%%%%%%%%%%%%%%%%%%%%%
